% ============================================================================
\chapter{Java~: le langage, la JVM, les bases}
\label{ch:java-bases}
% ============================================================================

\begin{objectif}
\begin{itemize}
  \item Installer un JDK, comprendre ce que sont JVM, bytecode, GC, et le mod\`ele \og ~\'ecrire une fois,
        ex\'ecuter partout~ \fg{}.
  \item Compiler et ex\'ecuter un programme, avec et sans outil de build.
  \item Ma\^itriser le typage statique, les primitifs, les conversions et la nullit\'e.
  \item \'Ecrire du Java moderne~: \verb|var|, enregistrements, expressions \verb|switch|, blocs de texte.
  \item Savoir lire une stack trace et un message du compilateur sans paniquer.
\end{itemize}
\end{objectif}

\section{La plateforme, en cinq mots}

\begin{description}
  \item[JDK] l'outil de d\'eveloppement~: compilateur (\verb|javac|), lanceur (\verb|java|),
        \verb|jshell|, \verb|javadoc|, \verb|jar|, analyeurs (\verb|jcmd|, \verb|jstack|, \verb|jmap|).
  \item[JVM] la machine virtuelle qui ex\'ecute le \emph{bytecode} (\fichier{.class}), just-\`a-temps
        (JIT), avec ramasse-miettes, et un \'ecosyst\`eme d'outils de profilage.
  \item[Version] une nouvelle version \textbf{tous les six mois}~; des versions \textbf{LTS}
        (17, 21, \emph{25}) soutenues des ann\'ees. En entreprise~: 17 ou 21 dominent.
        En 2026, \textbf{Java~21} est le choix par d\'efaut (virtual threads, records, patterns,
        s\'equenc\'e), \textbf{Java~25} apporte les versions de preview stabilis\'ees (structured
        concurrency, primitive types en preview, etc.).
  \item[Distributions] les binaires OpenJDK~: \link{Temurin (Adoptium)}{https://adoptium.net/fr/}
        (le choix par d\'efaut), \link{Amazon~Corretto}{https://aws.amazon.com/corretto/},
        \link{Azul~Zulu}{https://www.azul.com/downloads/}, \link{BellSoft~Liberica}{https://bell-sw.com/pages/downloads/},
        \link{GraalVM}{https://www.graalvm.org/} (compilation native, polyglotte).
  \item[Sp\'ecification] \link{JLS}{https://docs.oracle.com/javase/specs/} et les
        \link{JEP}{https://openjdk.org/jeps/0}--- chaque nouveaut\'es y est \'ecrite noir sur blanc.
\end{description}

\begin{terminal}
# Linux (Debian/Ubuntu) : le depot Temurin
sudo apt install -y temurin-21-jdk
# macOS
brew install --cask temurin@21
# Windows : https://adoptium.net/temurin/releases/?os=windows (MSI : cocher JAVA_HOME + PATH)
# Ou, multi-versions (recommande) :
brew install sdkman-cli && sdk install java 21.0.5-tem && sdk list java     # Linux/macOS
# Windows : scoop bucket add java && scoop install temurin21-jdk
# Verification :
java -version && javac -version && echo $JAVA_HOME
which java && readlink -f $(which java)     # ne pas etre surpris : /usr/bin/java est un lien alternatives
\end{terminal}

\subsection{Compiler, ex\'ecuter, et le \og ~single-file source launcher~ \fg{}}

\begin{codejava}
// HelloWorld.java --- le nom du fichier doit correspondre a la classe publique
public class HelloWorld {
    public static void main(String[] args) {
        System.out.println("Bonjour, Java !");
    }
}
\end{codejava}

\begin{terminal}
# Facon classique : compiler puis lancer
javac HelloWorld.java            # produit HelloWorld.class (le bytecode)
java HelloWorld                  # la JVM execute ; pas de ".class"

# Depuis Java 11 : lancer le fichier source directement (scripts, apprentissage)
java HelloWorld.java
# Avec arguments, proprietes, et classpath :
java -Xmx256m -Dapp.env=dev -cp lib/\*:target/classes com.exemple.App --port 8080
# JAR executable :
jar cfe app.jar com.exemple.App -C target/classes . && java -jar app.jar
# Decouvrir la JVM :
java --list-modules ; java -XX:+PrintFlagsFinal -version | grep -i heapsize ; jcmd VM.info
\end{terminal}

\begin{note}
Java est compil\'e \emph{et} ex\'ecut\'e par une JVM qui compile \`a chaud (C1 puis C2 JIT)~: le
d\'emarrage d'une application Java un peu grosse est lent, puis elle devient plus rapide que du JS
ou du PHP sur un calcul long. C'est pour cela qu'on utilise des cadres avec \og ~compilation
native~ \fg{} (GraalVM) ou des \emph{app/crac} pour le serverless~; et que dans ce livre, les mesures
de performance Java se font apr\`es \'echauffement (\link{JMH}{https://github.com/openjdk/jmh}).
\end{note}

\section{Types primitifs, variables et conversions}

\begin{codejava}
// Primitives : 8, pas d'objet, pas de null, empilees dans la pile d'appel.
boolean vif = true;
byte octet = 127;                    // -128..127 (rare, utile pour les protocols)
short petit = 32_000;                // -32 768..32 767
int n = 42;                          // -2^31..2^31-1 : c'est LE type entier par defaut
long grand = 9_000_000_000L;         // 64 bits : pour les id, les tailles, les timestamps
float virgule = 1.5f;                // 32 bits : ne l'utilise presque jamais
double d = 3.14159;                  // 64 bits IEEE-754 : le par defaut pour le flottant
char c = 'A';                        // 16 bits UTF-16 ; 0..65535 (les emoji debordent)

// 0b1010 = 10 ; 0x1F = 31 ; 1_000_000 s'ecrit avec des underscores ; 0777 est OCTAL (piege).
var inference = 42;                  // Java 10+ : le type est DEDUIT (ici int), pas " dynamique"
final int CONSTANTE = 7;             // non reassignable ; pour une veritable constante :
public static final double TVA = 0.20;

// Conversions
long l = n;                          // elargissement : implicite, sur
int revenu = (int) grand;             // retrecissement : EXPLICITE obligatoire (sinon erreur)
int arrondi = (int) Math.round(d);
long precise = Math.roundExact(2.5); // version 19+ : leve au lieu de silencieusement tronquer

// Division entiere : le piege n1 de Java
double demi = 1 / 2;                 // 0.0  <-- la division se fait en int AVANT conversion
double vraiDemi = 1.0 / 2;           // 0.5
double pourcentage = 5.0 * 2 / 100;  // 0.1 (et non 0)
int reste = 7 % -2;                   // 1 ; signe = celui du dividende

double infini = 1.0 / 0;              // Infinity : pas d'exception en flottant !
double nan = 0.0 / 0;                 // NaN
System.out.println(Double.isNaN(nan));          // true
System.out.println(nan == Double.NaN);          // false : NaN != NaN, toujours
System.out.println(0.1 + 0.2 == 0.3);           // false
System.out.println(Math.abs(0.1 + 0.2 - 0.3) < 1e-9);  // true
// En monetaire : JAMAIS double. Utilise BigDecimal (exact) ou des centimes en long.
var montant = java.math.BigDecimal.valueOf(12.50).multiply(java.math.BigDecimal.valueOf(3));
var arrondiUp = montant.setScale(2, java.math.RoundingMode.HALF_UP);   // 37.50

// Strings
String s = "Bonjour";
String bloc = """
    SELECT id, titre
    FROM   articles
    WHERE  statut = 'publie'
    """;                               // Java 15+ : text block, sans concatenations moches
var concat = "a" + 1 + true;           // "a1true" : evaluation gauche a droite
var interpo = "Resultat : " + arrondiUp.stripTrailingZeros().toPlainString();
System.out.printf("%-10s | %6.2f | %05d%n", "libelle", 12.34, 7);  // n = saut de ligne
String[] mots = bloc.lines().toList().toArray(String[]::new);
\end{codejava}

\subsection{Nullit\'e, types enveloppes, et \texttt{Optional}}

\begin{codejava}
Integer boxee = null;                  // les wrappers (Integer, Long, Double...) peuvent etre null
int casse = boxee;                      // NullPointerException a l'unboxing : le classique
long somme = liste.stream().mapToLong(Moteur::getDuree).sum();  // si un getDuree est Long null -> NPE

// Comparaison : equals/toujours pour les objets, == seulement pour les primitives
Integer a = 128, b = 128;
System.out.println(a == b);             // FALSE  (cache des Integer uniquement -128..127 !)
System.out.println(Objects.equals(a, b));   // true : LE reflexe a prendre
System.out.println("x".equals(s));         // true ; "x" == s : comparaison de reference -> faux

// Optional : un conteneur de valeur absente, pas un champ de classe
public Optional<Article> trouver(long id) {
    return Optional.ofNullable(repo.sel(id));
}
Optional<Article> o = trouver(1);
o.ifPresent(a -> System.out.println(a.titre()));
String titre = o.map(Article::titre).orElse("sans titre");
Article article = o.orElseThrow(() -> new Introuvable("id 1"));
o.filter(a -> a.publie()).ifPresent(this::notifier);
// Anti-patterns : o.get() sans isPresent(), Optional comme parametre ou champ, .orElse(calculCouteux())
// (le calcul est evalue dans tous les cas : orElseGet le rend paresseux)
\end{codejava}

\section{Op\'erateurs, conditions et boucles}

\begin{codejava}
// Arithmetique et incrementation : comme en JS/C, aux differences notees
int i = 0; i++; ++i; i += 2; i <<= 3; i >>>= 1;         // >>> : decale logique (sans signe)
boolean b = !b;                                         // pas de surcharge operateur en Java
long hachage = "cle".hashCode();                        // pas de `+` sur des types arbitraires

// Comparaison a trois voies : le spaceship de Java
Integer.compare(x, y);                                  // -1/0/1 (utilise dans les comparators)
Long.compare(a, b);
"b".compareTo("a");
Comparator.comparing(Article::date).thenComparing(Article::titre).reversed();

// Conditions
if (a > 0 && b < 10) { } else if (a == 0) { } else { }
String libelle = switch (note) {                        // Java 14+ : switch EXPRESSION
    case 20 -> "parfait";
    case 16, 17, 18, 19 -> "tres bien";
    case 10, 11, 12, 13, 14, 15 -> "bien";
    default -> {
        String r = note < 0 ? "impossible" : "a revoir";
        yield r;                                        // yield = "return" dans un switch
    }
};
// switch sur des types et des motifs (Java 21) : le pattern matching
String decrire(Object o) {
    return switch (o) {
        case Integer n when n > 100 -> "grand entier " + n;
        case Integer n -> "petit entier " + n;
        case String s when s.isBlank() -> "vide";
        case String s -> "texte de " + s.length() + " caracteres";
        case int[] tab -> "tableau de " + tab.length + " entiers";
        case null -> "rien";                             // Java 21 : le cas null est autorise
        default -> "autre chose";
    };
}

// Boucles
for (int j = 0; j < 10; j++) { }
for (var mot : "un deux".split(" ")) { }
for (int k = 0, m = 10; k < m; k++, m--) { }
while (it.hasNext()) { }
do { travail(); } while (tropTot());
// Etiquettes : a connaitre pour lire du code, a eviter d'ecrire
externe: for (int x = 0; x < 10; x++) { for (int y = 0; y < 10; y++) { if (y > x) continue externe; } }
\end{codejava}

\section{Tableaux~: le type brut dont tout le monde h\'erite dans sa t\^ete}

\begin{codejava}
int[] notes = new int[5];                 // taille fixe, initialisee a 0
int[] litteral = {14, 9, 17, 12};
String[] villes = { "Dakar", "Thies", "Saint-Louis" };
notes.length;                              // une PROPRIETE, pas une methode
villes[0] = "Thi\`es";
//ville[5];                                 // ArrayIndexOutOfBoundsException (jamais de lecture floue)

var matrice = new int[3][4];               // tableaux imbriques (lignes de tailles differentes OK)
matrice[0][1] = 7;

Arrays.sort(litteral);
Arrays.copyOf(litteral, 10);
Arrays.equals(litteral, copie);
Arrays.fill(notes, -1);
Arrays.mismatch(a, b);                      // Java 11+ : index de difference, -1 si egales
Arrays.asList("a", "b");                    // vue FIXE sur le tableau -- pas un ArrayList !
List.of(1, 2, 3);                            // liste immuable (Java 9+)
new ArrayList<>(List.of(1, 2, 3));           // la version modifiable

// Impression (jamais System.out.println(tab) : tu obtiens [I@1a2b3c)
System.out.println(Arrays.toString(litteral));
System.out.println(Arrays.deepToString(matrice));
System.out.printf("notes = %s%n", Arrays.toString(notes));

// Copie rapide et allocation : les deux utilitaires bas niveau
System.arraycopy(src, 0, dest, 10, n);     // ou : Arrays.copyOfRange(src, debut, fin)
var montee = new java.util.RandomGenerator.JdkRandom(1).ints(1000, 0, 50).toArray();
// Tableaux et genericite : impossible de creer un List<String>[] typable -> on utilise des Collections.
\end{codejava}

\section{M\'ethodes, arguments, et la surcharge}

\begin{codejava}
public final class Calculs {
    private Calculs() { }                        // classe utilitaire : pas d'instance

    // Surcharge : meme nom, signatures differentes (les types, pas les noms de parametres)
    public static int max(int a, int b) { return Math.max(a, b); }
    public static double max(double a, double b) { return Math.max(a, b); }
    public static int max(int... valeurs) {       // varargs : un tableau en entree
        if (valeurs.length == 0) { throw new IllegalArgumentException("au moins une valeur"); }
        int m = valeurs[0];
        for (var v : valeurs) { if (v > m) { m = v; } }
        return m;
    }
    public static void melanger(List<String> l) { java.util.Collections.shuffle(l); }
}

// Pas de passage par reference : les objets passent par copie de la reference
static void tente(Article a) { a = new Article("autre"); }   // l'appelant ne le voit pas
static void marche(Article a) { a.marquerVu(); }              // mais muter l'objet, si

// En Java, une methode renvoie UN seul resultat : pour en renvoyer plusieurs -> un record
public record DivResultat(int quotient, int reste) { }
public static DivResultat diviser(int a, int b) { return new DivResultat(a / b, a % b); }

// Recursion : pas de TCO garantie par la JVM -> sur de tres grandes profondeurs, une boucle.
static long fib(int n, long a, long b) { return n == 0 ? a : fib(n - 1, b, a + b); }

// Memo : les methodes d'une classe ne sont jamais globales ; tout est statique ou membre.
\end{codejava}

\section{Lire les messages~: le compilateur et la pile d'appels}

\begin{codejava}
/*
  Exemples reellement frequents, avec leur traduction :

  error: cannot find symbol                -> classe non importee, faute de frappe, ou methode
                                             inexistante ; verifie la casse et le package
  error: variable x might not have been initialized -> champ local non initialise (en Java, un
                                             local n'a AUCUNE valeur par defaut)
  error: incompatible types: int cannot be converted to String -> String.valueOf(n) ou "" + n
  error: unreported exception IOException  -> les exceptions verifiees : try/catch ou throws
  error: method does not override or implement a method from a supertype
                                          -> signature legerement differente : mets @Override
  Exception in thread "main" java.lang.NullPointerException: Cannot invoke "String.length()"
       because "s" is null                 -> depuis Java 14, le message dit QUI est null. Lis-le.
  java.lang.ArrayIndexOutOfBoundsException: Index 5 out of bounds for length 5
  java.util.ConcurrentModificationException -> tu modifies une collection pendant son iteration
  java.lang.StackOverflowError              -> recursion sans cas de base (ou cycle d'objets)
*/
\end{codejava}

\begin{exercice}
\begin{enumerate}
  \item \'Ecris \verb|FizzBuzz|, \verb|estPremier|, \verb|pgcd| (Euclide, puis version binaire),
        \verb|converser(int[])| (sur place), \verb|moyenneMobile(double[], int)|~: tous en
        m\'ethodes statiques d'une classe utilitaire priv\'ee.
  \item Pr\'edis puis v\'erifie~: \verb|1/2|, \verb|1.0/2|, \verb|7\%3|, \verb|-7\%3|,
        \verb|(int) 3.99|, \verb|Math.round(-2.5)|, \verb|Double.doubleToLongBits(0.0)|,
        \verb|Integer.MAX\_VALUE + 1|, \verb|0.1 + 0.2 == 0.3|, \verb|"1" + 2|, \verb|2 + "1"|.
  \item Convertis un nombre en romains (\verb|1..3999|), puis l'inverse~; puis un nombre en
        base quelconque (2 \`a 36) et reviens, en utilisant \verb|Integer.toString(n, base)|
        et \verb|Integer.parseInt(s, base)|~: compare \`a ta version manuscrite.
  \item \'Ecris \verb|Textes|~: \verb|inverser|, \verb|estPalindrome| (en ignorant ponctuation et
        accents via \verb|java.text.Normalizer|), \verb|compterMots|, \verb|capitaliser|,
        \verb|estAnagramme|, \verb|slug|~: en utilisant \verb|StringBuilder| et les m\'ethodes
        \verb|chars()|/ \verb|codePoints()|.
  \item \'Ecris une classe \verb|Matrice| (tableau 2D) avec \verb|ajoute|, \verb|multiplier|,
        \verb|transposer|, \verb|toString| et \verb|equals|/ \verb|hashCode| corrects.
  \item Ouvre \verb|jshell|~: teste chaque snippet de ce chapitre une ligne \`a la fois.
        C'est le meilleur moyen de v\'erifier une hypoth\`ese sans cr\'eer de projet.
\end{enumerate}
\end{exercice}

\begin{ressources}
\begin{itemize}
  \item Dev.java (la porte officielle moderne) --- \url{https://dev.java/learn/}
  \item Documentations JDK (javadoc de r\'ef\'erence) --- \url{https://docs.oracle.com/en/java/javase/21/docs/api/index.html}
  \item \link{Baeldung}{https://www.baeldung.com/java-tutorial} (recettes concises, enorme index)
  \item \emph{Introduction to Java Programming} -- \url{https://java-programming.baoyj.com/} (libre)
  \item \link{JEP index}{https://hiphopjava.github.io/jep-in-Java-FR/} (JEP traduits en francais)
\end{itemize}
\end{ressources}

\index{Java!bases}\index{JDK}\index{JVM}\index{primitives}\index{Optional}
